\documentclass[11pt,a4paper]{article}

\usepackage[T1]{fontenc}
\usepackage{lmodern}
\usepackage{amsmath,amssymb,amsthm,mathtools}
\usepackage[margin=27mm]{geometry}
\usepackage{microtype}
\usepackage[hidelinks]{hyperref}

\newcommand{\R}{\mathbb R}
\newcommand{\Z}{\mathbb Z}
\newcommand{\ind}{\mathbf 1}
\DeclareMathOperator{\supp}{supp}
\DeclareMathOperator{\dist}{dist}
\newtheorem*{targettheorem}{Target theorem}
\newtheorem{lemma}{Lemma}[section]
\newtheorem{proposition}[lemma]{Proposition}

\setlength{\emergencystretch}{2em}
\allowdisplaybreaks

\title{Bilinear smoothing along distinct monomials\\[4pt]
\large A formalization blueprint using the Christ--Durcik--Roos method}
\author{}
\date{}

\begin{document}
\maketitle

\begin{targettheorem}
Let $\alpha_1,\alpha_2$ be distinct positive integers and let
$0<a<b<\infty$. There exist constants
\[
C=C(a,b,\alpha_1,\alpha_2)<\infty,
\qquad \delta>0,
\]
such that, for every $\psi\in C^1(\R)$ supported in $[a,b]$, every
$\lambda\geq1$, and every $g_1,g_2\in\mathcal S(\R^2)$ satisfying
\[
\supp\widehat g_j\subseteq\{\xi\in\R^2:|\xi_j|\geq\lambda\}
\quad\text{for at least one }j\in\{1,2\},
\]
the operator
\[
\mathcal A_\psi(g_1,g_2)(x,y)
=\int_{\R}\psi(t)g_1(x+t^{\alpha_1},y)
                    g_2(x,y+t^{\alpha_2})\,dt
\]
satisfies
\[
\|\mathcal A_\psi(g_1,g_2)\|_1
\leq C\bigl(\|\psi\|_\infty+\|\psi'\|_1\bigr)
\lambda^{-\delta}\|g_1\|_2\|g_2\|_2.
\tag{T}\label{eq:target}
\]
The exponent can be chosen independently of $\psi$. The proof below gives,
without attempting optimization, the universal choice
\[
\delta=\frac{1}{2\,935\,296}.
\]
\end{targettheorem}

The argument follows the Christ--Durcik--Roos method, with a modified flow.
The monomial extension is discussed in their paper~\cite{CDR}. Below, the
structural decomposition, sublevel estimate, local analytic estimate,
interpolation, globalization, and $C^1$-cutoff reduction are proved.
Neither the Christ--Durcik--Roos smoothing theorem nor a general
sublevel-set theorem is used as a black box.

\section{Framework and elementary Fourier facts}

We first work with a real number $r>1$, an interval
\[
I=[A,B]\subset(0,\infty),
\]
and
\[
Q(t)=t^r,\qquad q(t)=Q'(t)=rt^{r-1},\qquad
\Gamma(t)=(-t,Q(t)).
\]
On a fixed neighborhood of $I$, the functions $q$ and $q'$ have positive
lower bounds, and all derivatives of $Q$ have finite upper bounds.

Define
\[
T_\psi(f_1,f_2)(x,y)
=\int_I\psi(t)f_1(x+t,y)f_2(x,y+Q(t))\,dt.
\tag{1.1}\label{eq:normalized-operator}
\]
All unspecified constants below may depend on $r,I$, fixed auxiliary smooth
functions, and the size of a fixed spatial localization region. They do not
depend on the frequency parameter or the functions being estimated.

We use the Fourier transform
\[
\widehat f(\xi)=\int f(x)e^{-2\pi i x\cdot\xi}\,dx.
\]
The foundational results used are Fubini--Tonelli, H\"older and
Cauchy--Schwarz, convergence theorems for Lebesgue integrals, density in
$L^p$, Fourier inversion and Plancherel, elementary differential calculus,
the inverse function theorem, change of variables, and the maximum-modulus
principle. The $L^2$ Fourier theory in mathlib is documented
in~\cite{mathlib-fourier}; the maximum-modulus principle is documented
in~\cite{mathlib-maximum}.

\subsection{The elementary \texorpdfstring{$L^2\times L^2\to L^1$}{L2 x L2 to L1} bound}

For arbitrary $\psi\in L^1(I)$, translation invariance and Cauchy--Schwarz give
\[
\|T_\psi(f_1,f_2)\|_1
\leq\int_I|\psi(t)|
\|f_1(\cdot+(t,0))f_2(\cdot+(0,Q(t)))\|_1\,dt
\leq\|\psi\|_1\|f_1\|_2\|f_2\|_2.
\tag{1.2}\label{eq:trivial}
\]
Consequently, $T_\psi$ extends continuously to $L^2\times L^2$, and
subsequent $L^2$-approximation arguments are justified by
\eqref{eq:trivial}.

\subsection{Smooth frequency cutoffs}

If $\vartheta\in C_c^\infty(\R)$, the multiplier $\vartheta(\xi/L)$ has
convolution kernel
\[
K_L(s)=L\check\vartheta(Ls).
\]
Integration by parts in the inverse Fourier integral gives
\[
|K_L(s)|\leq C_NL(1+L|s|)^{-N}.
\tag{1.3}\label{eq:kernel-decay}
\]
Scaling therefore gives
\[
\|K_L\|_1\leq C,\qquad
\|K_L'\|_1\leq CL,\qquad
\int |s||K_L(s)|\,ds\leq CL^{-1}.
\tag{1.4}\label{eq:kernel-norms}
\]
These estimates also hold for convolution in just one coordinate of $\R^2$.

The bound
\[
\|K*f\|_p\leq\|K\|_1\|f\|_p,\qquad 1\leq p\leq\infty,
\tag{1.5}\label{eq:young-l1}
\]
follows, for finite $p$, by applying Jensen's inequality to the probability
measure $|K(s)|\,ds/\|K\|_1$, then integrating in $x$. The case $p=\infty$
is immediate; when $\|K\|_1=0$ the assertion is trivial.

\subsection{An autocorrelation identity}

For $f\in L^2(\R)$, put
\[
D_sf(x)=f(x+s)\overline{f(x)}
\]
and
\[
\mathcal U_R(f)
=\int_{\R}\int_{|\xi|\leq R}
|\widehat{D_sf}(\xi)|^2\,d\xi\,ds.
\]
Then
\[
\mathcal U_R(f)
=\iint_{|\xi-\eta|\leq R}
|\widehat f(\xi)|^2|\widehat f(\eta)|^2\,d\xi\,d\eta.
\tag{1.6}\label{eq:autocorrelation}
\]
In particular,
\[
\mathcal U_R(f)
\leq\|f\|_2^2
\sup_{\substack{J\subset\R\\ |J|=2R}}
\int_J|\widehat f(\xi)|^2\,d\xi.
\tag{1.7}\label{eq:autocorrelation-bound}
\]
For Schwartz functions, Fourier inversion gives
\[
\widehat{D_sf}(\xi)
=\int e^{2\pi i s\eta}
\widehat f(\eta)\overline{\widehat f(\eta-\xi)}\,d\eta.
\]
Plancherel in $s$, followed by a change of variables, proves
\eqref{eq:autocorrelation}.

For general $L^2$ functions, approximate by Schwartz functions. The required
continuity follows from
\[
\|D_sf-D_sg\|_{L^2(ds\,dx)}
\leq\|f-g\|_2\bigl(\|f\|_2+\|g\|_2\bigr),
\]
which is obtained by expanding the difference and applying Fubini.
Thus \eqref{eq:autocorrelation} passes to the limit.

We will also use the direct Fubini identity
\[
\int_{\R}\|D_sf\|_2^2\,ds=\|f\|_2^4.
\tag{1.8}\label{eq:autocorrelation-total}
\]

\section{The frequency decomposition}

This is the structural part of the Christ--Durcik--Roos argument. The
following construction proves exactly the properties needed later.

\begin{lemma}\label{lem:decomposition}
Let $R\geq1$, $0<\rho<1$, and
$f\in L^2(\R)\cap L^\infty(\R)$. There is a decomposition
\[
f=f_\sharp+f_\flat
\]
such that
\[
|\widehat{f_\sharp}|\leq|\widehat f|,
\qquad |\widehat{f_\flat}|\leq|\widehat f|,
\tag{2.1}\label{eq:decomp-domination}
\]
and
\[
\sup_{|J|=2R}\int_J|\widehat{f_\flat}|^2
\leq\rho\|f\|_2^2.
\tag{2.2}\label{eq:decomp-flat}
\]
Moreover,
\[
f_\sharp(x)=\sum_{\nu=1}^Nh_\nu(x)e^{2\pi i\omega_\nu x},
\qquad N\leq C\rho^{-1},
\tag{2.3}\label{eq:decomp-sharp}
\]
where
\[
\supp\widehat h_\nu\subseteq[-R,R],\qquad
\|\partial_x^kh_\nu\|_\infty\leq C_kR^k\|f\|_\infty.
\tag{2.4}\label{eq:decomp-amplitudes}
\]
Every nonzero summand can be chosen with
\[
\dist\bigl(\omega_\nu,\supp\widehat f\bigr)\leq R.
\tag{2.5}\label{eq:decomp-centers}
\]
\end{lemma}

\begin{proof}
Choose a nonnegative $\vartheta\in C_c^\infty((-1,1))$ satisfying
\[
\sum_{n\in\Z}\vartheta(u-n)=1.
\]
Such a function is obtained by dividing a nonnegative bump, positive on
$[-3/4,3/4]$, by the sum of its integer translates.

Write $E=\|f\|_2^2$. The case $E=0$ is immediate. Declare $n\in\Z$ to be
selected when
\[
\int_{(n-4)R}^{(n+4)R}|\widehat f(\xi)|^2\,d\xi>\rho E.
\]
Each point belongs to at most nine of these intervals, so there are at most
$9\rho^{-1}$ selected integers.

Define
\[
\widehat{f_\sharp}(\xi)
=\sum_{n\ \mathrm{selected}}
\vartheta(\xi/R-n)\widehat f(\xi),
\]
and set $f_\flat=f-f_\sharp$. The multipliers defining the two parts lie
between zero and one, proving \eqref{eq:decomp-domination}.

Let $J$ have length $2R$. If an unselected multiplier is nonzero somewhere
on $J$, then $J$ is contained in its corresponding interval
$[(n-4)R,(n+4)R]$. Hence
\[
\int_J|\widehat{f_\flat}|^2\leq\rho E.
\]
If no unselected multiplier is nonzero on $J$, the integral is zero. This
proves \eqref{eq:decomp-flat}.

For a selected $n$, define
\[
h_n(x)=e^{-2\pi i nRx}\bigl[\vartheta(D/R-n)f\bigr](x).
\]
Its Fourier support lies in $[-R,R]$. Its convolution representation,
followed by differentiation of the fixed Schwartz kernel, proves
\eqref{eq:decomp-amplitudes}. Discarding zero summands gives
\eqref{eq:decomp-centers}.
\end{proof}

Two consequences require explicit attention. First,
\eqref{eq:autocorrelation-bound} and \eqref{eq:decomp-flat} give
\[
\mathcal U_R(f_\flat)\leq\rho\|f\|_2^4.
\tag{2.6}\label{eq:flat-correlation}
\]
Second, suppose $\chi$ is a spatial cutoff whose Fourier transform has
$L^1$ norm at most $C$. For every interval $J$ of length $2R$, the
convolution formula and the triangle inequality in $L^2(J)$ give
\[
\|\ind_J\widehat{\chi f_\flat}\|_2
\leq\int |\widehat\chi(\eta)|
\|\ind_{J-\eta}\widehat{f_\flat}\|_2\,d\eta
\leq C\rho^{1/2}\|f\|_2.
\]
If also $|\chi|\leq1$, then
\[
\mathcal U_R(\chi f_\flat)\leq C\rho\|f\|_2^4.
\tag{2.7}\label{eq:localized-flat-correlation}
\]
The constant is unchanged by translating or dilating a fixed cutoff.

The construction can be performed measurably on fibers of a function on
$\R^2$: all selection conditions are measurable integrals. Enumerate the
selected integers using the fixed ordering $0,1,-1,2,-2,\ldots$, and pad
the expansion with zero summands. This produces measurable frequency
functions in the passive coordinate.

\section{The modified flow and the sublevel estimate}

\begin{lemma}\label{lem:sublevel}
Fix a bounded spatial box $D\subset\R^2$ and $c_0>0$. Let
$\alpha,\beta:\R^2\to\R$ be measurable. Suppose either
\[
|\alpha(z)|\geq c_0\quad\text{for all }z,
\]
or
\[
|\beta(z)|\geq c_0\quad\text{for all }z.
\]
Then, for $0<\varepsilon\leq1$,
\[
\left|\left\{(x,y,t)\in D\times I:
|\alpha(x+t,y)-q(t)\beta(x,y+Q(t))|\leq\varepsilon
\right\}\right|
\leq C\varepsilon^\kappa,
\qquad\kappa=\frac1{39}.
\tag{3.1}\label{eq:sublevel}
\]
\end{lemma}

We prove this in three steps.

\subsection{The vector field and its quantitative nondegeneracy}

For $\mathbf t=(t_1,t_2,t_3)$, define
\[
\theta_1(\mathbf t)=-t_1+t_2-t_3,
\]
\[
\theta_2(\mathbf t)=Q(t_1)-Q(t_2)+Q(t_3),
\]
and
\[
\nu(\mathbf t)=\frac{q(t_2)}{q(t_1)q(t_3)}.
\tag{3.2}\label{eq:flow-coordinates}
\]
The appropriate vector field is
\[
V=\nabla\theta_1\times\nabla\theta_2.
\]
Explicitly,
\[
V=\bigl(q(t_3)-q(t_2)\bigr)\partial_{t_1}
 +\bigl(q(t_3)-q(t_1)\bigr)\partial_{t_2}
 +\bigl(q(t_2)-q(t_1)\bigr)\partial_{t_3}.
\tag{3.3}\label{eq:vector-field}
\]
By construction,
\[
V\theta_1=V\theta_2=0.
\tag{3.4}\label{eq:flow-invariants}
\]
Because $q(t)=rt^{r-1}$,
\[
V\nu=(r-1)\nu\left(
\frac{q(t_2)-q(t_3)}{t_1}
+\frac{q(t_3)-q(t_1)}{t_2}
+\frac{q(t_1)-q(t_2)}{t_3}\right).
\tag{3.5}\label{eq:flow-derivative}
\]
We claim that
\[
|V\nu(\mathbf t)|\geq
c\prod_{1\leq i<j\leq3}|t_i-t_j|,
\qquad\mathbf t\in I^3.
\tag{3.6}\label{eq:flow-nondegeneracy}
\]
To prove it, set $p=r-1>0$ and $u_i=t_i^p$. Up to sign, the parenthesis
in \eqref{eq:flow-derivative} is the determinant with rows
\[
(1,q(t_i),t_i^{-1}),\qquad i=1,2,3.
\]
Since $q(t_i)=ru_i$ and $t_i^{-1}=u_i^{-1/p}$, its absolute value is
\[
r\left|\det\bigl[(1,u_i,u_i^{-1/p})\bigr]_{i=1}^3\right|.
\]
For distinct $u_i$, let $P$ be the quadratic polynomial interpolating
$F(u)=u^{-1/p}$ at the three points. Two applications of Rolle's theorem
show that the leading coefficient of $P$ is $F''(u_*)/2$ for some
intermediate $u_*$. Since
\[
F''(u)=\frac1p\left(\frac1p+1\right)u^{-1/p-2}>0,
\]
this coefficient is bounded below on the relevant compact interval.
Subtracting the constant and linear parts of $P$ in the determinant gives
\[
\left|\det\bigl[(1,u_i,F(u_i))\bigr]_{i=1}^3\right|
\geq c\prod_{i<j}|u_i-u_j|.
\]
The derivative of $t\mapsto t^p$ has a positive lower bound on $I$, so
\[
|u_i-u_j|\geq c|t_i-t_j|.
\]
Finally, $\nu$ has a positive lower bound on $I^3$. This proves
\eqref{eq:flow-nondegeneracy}, including coincident points, where the
right-hand side is zero.

Thus the map
\[
H(\mathbf t)=\bigl(\theta_1(\mathbf t),\theta_2(\mathbf t),\nu(\mathbf t)\bigr)
\tag{3.7}\label{eq:flow-map}
\]
satisfies
\[
|\det DH(\mathbf t)|=|V\nu(\mathbf t)|
\geq c\prod_{i<j}|t_i-t_j|.
\tag{3.8}\label{eq:jacobian-bound}
\]
This is the flow calculation needed for the monomial extension. Where
$V\nu\neq0$, the time-rescaled field
\[
W=\frac{V}{V\nu}
\]
satisfies
\[
W\theta_1=W\theta_2=0,\qquad W\nu=1.
\]
Consequently, in the coordinates $H$, its local flow is simply translation
in the third coordinate. No explicit solution of the original
differential equation is required.

\subsection{A sublevel estimate in the flow coordinates}

For every real number $a_0$ with $|a_0|\geq c_0$, every measurable
$b:\R^2\to\R$, and every $z_0\in\R^2$, we have
\[
\left|\left\{\mathbf t\in I^3:
|a_0\nu(\mathbf t)-b(z_0+(\theta_1(\mathbf t),\theta_2(\mathbf t)))|
\leq\varepsilon\right\}\right|
\leq C\varepsilon^{1/13}.
\tag{3.9}\label{eq:coordinate-sublevel}
\]
Here are the quantitative details.

Fix $0<\eta<1$. The set where some $|t_i-t_j|<\eta$ has measure at most
$C\eta$. On its complement,
\[
|\det DH|\geq c\eta^3.
\tag{3.10}\label{eq:good-jacobian}
\]
The first and second derivatives of $H$ are uniformly bounded on a fixed
neighborhood of $I^3$. Cover $I^3$ by cubes of side length $c_1\eta^3$,
with $c_1>0$ sufficiently small. There are at most $C\eta^{-9}$ such cubes.

Consider a cube meeting the complement of the exceptional set, and choose
such a point $\mathbf t_0$ in it. The cofactor formula and
\eqref{eq:good-jacobian} give
\[
\|(DH(\mathbf t_0))^{-1}\|\leq C\eta^{-3}.
\]
By the bound on the second derivatives and the choice of $c_1$,
\[
\sup_{\mathbf t\text{ in the cube}}
\|DH(\mathbf t)-DH(\mathbf t_0)\|
\leq\frac12\|(DH(\mathbf t_0))^{-1}\|^{-1}.
\]
Integrating derivatives along line segments shows that $H$ is injective
on the cube.

For a measurable subset $E$ of its good part, change of variables
therefore gives
\[
|E|\leq C\eta^{-3}|H(E)|.
\]
The image $H(I^3)$ lies in a fixed bounded box. In the image coordinates
$(u,v)\in\R^2\times\R$, the sublevel condition is
\[
|a_0v-b(z_0+u)|\leq\varepsilon.
\]
For every $u$, the permitted $v$-set has length at most
$2\varepsilon/c_0$. Fubini gives image measure at most $C\varepsilon$.

Summing over the cubes yields
\[
|\text{sublevel set}|\leq C\eta+C\varepsilon\eta^{-12}.
\]
Taking $\eta=\varepsilon^{1/13}$ proves
\eqref{eq:coordinate-sublevel} (the endpoint $\varepsilon=1$ follows by
increasing the constant).

Measurable functions can first be replaced by Borel representatives. The
changes of variables above then establish all required measurability
statements; changes on null sets do not affect the original
three-dimensional sublevel set.

\subsection{Three successive relations}

After the measure-preserving substitution $z=(x+t,y)$, it suffices to
estimate
\[
E=\{(z,t)\in D_0\times I:
|\alpha(z)-q(t)\beta(z+\Gamma(t))|\leq\varepsilon\},
\tag{3.11}\label{eq:incidence-set}
\]
where $D_0$ is a fixed bounded box.

Assume first that $|\alpha|\geq c_0$. Write $\mu=|E|$, and suppose
$\mu>0$. Define
\[
d_A(z)=\int_I\ind_E(z,t)\,dt,\qquad
 d_B(w)=\int_I\ind_E(w-\Gamma(t),t)\,dt.
\]
Their supports lie in fixed bounded boxes of measures $M_A,M_B$.

The portion of $E$ where
\[
d_A(z)<\frac{\mu}{4M_A}
\]
has measure at most $\mu/4$. The portion where
\[
d_B(z+\Gamma(t))<\frac{\mu}{4M_B}
\]
also has measure at most $\mu/4$, by translation in $z$. Hence
\[
\int_Ed_A(z)d_B(z+\Gamma(t))\,dz\,dt\geq c\mu^3.
\tag{3.12}\label{eq:three-edge-count}
\]
For $z_0\in D_0$, let $\mathcal A(z_0)\subset I^3$ consist of the triples
for which
\[
(z_0,t_1)\in E,
\]
\[
(z_0+\Gamma(t_1)-\Gamma(t_2),t_2)\in E,
\]
and
\[
(z_0+\Gamma(t_1)-\Gamma(t_2),t_3)\in E.
\]
Fubini and the substitution $z=z_0+\Gamma(t_1)-\Gamma(t_2)$ give the exact
identity
\[
\int_{D_0}|\mathcal A(z_0)|\,dz_0
=\int_Ed_A(z)d_B(z+\Gamma(t))\,dz\,dt.
\]
Thus some $z_0$ satisfies
\[
|\mathcal A(z_0)|\geq c\mu^3.
\tag{3.13}\label{eq:large-fiber}
\]
For a triple in $\mathcal A(z_0)$, the three defining inequalities, and the
upper and lower bounds on $q$, imply
\[
\left|\alpha(z_0)\frac{q(t_2)}{q(t_1)q(t_3)}
-\beta\bigl(z_0+\Gamma(t_1)-\Gamma(t_2)+\Gamma(t_3)\bigr)\right|
\leq C\varepsilon.
\tag{3.14}\label{eq:three-relations}
\]
Indeed, solve the first relation for the first value of $\beta$, substitute
in the second, and then substitute in the third. All three error
coefficients are bounded because $q$ is bounded above and below.

Applying \eqref{eq:coordinate-sublevel} to \eqref{eq:three-relations} gives
\[
\mu^3\leq C\varepsilon^{1/13}.
\]
If $C\varepsilon>1$, the same inequality follows from the bounded volume
of $I^3$, after adjusting the constant. This proves \eqref{eq:sublevel}
with $\kappa=1/39$ when $|\alpha|\geq c_0$.

If instead $|\beta|\geq c_0$, change variables $w=z+\Gamma(t)$ and divide
the defining inequality by $q(t)$. This interchanges $\alpha,\beta$,
replaces $\Gamma$ by $-\Gamma$, and replaces $q$ by $q^{-1}$. The new
alternating-coordinate map is $-(\theta_1,\theta_2)$, and the new scalar
coordinate is $\nu^{-1}$. Its corresponding derivative is
\[
V(\nu^{-1})=-\nu^{-2}V\nu,
\]
so the same determinant lower bound and proof apply. This completes the
proof of Lemma~\ref{lem:sublevel}.\qed

\section{A positive \texorpdfstring{$L^p$}{Lp}-improving bound}

We next prove the auxiliary estimate
\[
\|T_\psi(f_1,f_2)\|_1
\leq C\|\psi\|_\infty
\|f_1\|_{12/7}\|f_2\|_{12/7}.
\tag{4.1}\label{eq:improving}
\]
Let $\mu$ be the measure on $\Gamma(I)$ obtained by pushing forward $dt$,
and let $\widetilde\mu$ be its reflection.

The measure $\mu*\widetilde\mu$ is the pushforward of $ds\,dt$ under
\[
(s,t)\longmapsto\Gamma(s)-\Gamma(t).
\]
Off the diagonal, this map is injective. Indeed, its first coordinate
fixes $t-s$; for fixed nonzero $t-s$, the second coordinate is strictly
monotone in $s$, since $q$ is strictly increasing.

Its Jacobian has absolute value
\[
|q(t)-q(s)|\asymp|t-s|.
\]
Consequently, $\mu*\widetilde\mu$ has a density $k$, and change of variables
gives
\[
\int_{\R^2}|k(w)|^{3/2}\,dw
=\int_{I^2}|q(t)-q(s)|^{-1/2}\,ds\,dt<\infty.
\tag{4.2}\label{eq:difference-density}
\]
We use the convolution inequality
\[
\|u*v\|_3\leq\|u\|_{3/2}\|v\|_{3/2}.
\tag{4.3}\label{eq:young-special}
\]
For completeness, H\"older with three exponents equal to three, applied to
\[
|u(x-y)v(y)|
=\bigl(|u(x-y)v(y)|^{1/2}\bigr)
|u(x-y)|^{1/2}|v(y)|^{1/2},
\]
gives
\[
|u*v(x)|^3
\leq\bigl(|u|^{3/2}*|v|^{3/2}\bigr)(x)
\|u\|_{3/2}^{3/2}\|v\|_{3/2}^{3/2}.
\]
Integration proves \eqref{eq:young-special}.

It follows that
\[
\|\mu*f\|_2^2=\langle k*f,f\rangle
\leq\|k*f\|_3\|f\|_{3/2}
\leq C\|f\|_{3/2}^2.
\tag{4.4}\label{eq:curve-convolution}
\]
The same holds for $\widetilde\mu$.

For nonnegative $F,G$, write
\[
B(F,G)=\int_{\R^2}\int_IF(z)G(z+\Gamma(t))\,dt\,dz.
\]
Equation~\eqref{eq:curve-convolution} and its adjoint give
\[
B(F,G)\leq C\|F\|_{3/2}\|G\|_2,\qquad
B(F,G)\leq C\|F\|_2\|G\|_{3/2}.
\]
Cauchy--Schwarz in $(z,t)$, using
\[
F(z)G(z+\Gamma(t))
=\bigl(F(z)^{8/7}G(z+\Gamma(t))^{6/7}\bigr)^{1/2}
 \bigl(F(z)^{6/7}G(z+\Gamma(t))^{8/7}\bigr)^{1/2},
\]
therefore gives
\[
B(F,G)\leq C\|F\|_{12/7}\|G\|_{12/7}.
\]
Finally, the substitution $z=(x+t,y)$ in the absolute-value bound for
$T_\psi$ proves \eqref{eq:improving}.

\section{The local correlation estimate}

Fix $\chi\in C_c^\infty(\R^2)$, and write
\[
T_{\mathrm{loc}}(f_1,f_2)=\chi T_\psi(f_1,f_2).
\]
For this section and the next, assume
\[
\|\psi\|_{C^{20}}:=\max_{0\leq k\leq20}\|\psi^{(k)}\|_\infty\leq1.
\tag{5.1}\label{eq:smooth-cutoff-normalization}
\]

\begin{lemma}\label{lem:local-correlation}
Let $0<h\leq1$, $L\geq h^{-1}$, $K\geq1$, and $0<\rho<1$. For $j=1,2$,
suppose that $v_{j,m}$, $m\in\Z^2$, is supported in a cube of side $Ch$
centered at $hm$, and
\[
\|v_{j,m}\|_\infty\leq B_j,\qquad
\|\partial_jv_{j,m}\|_\infty\leq B_jL.
\tag{5.2}\label{eq:local-derivatives}
\]
Assume also
\[
\int_{\R}\|v_{j,m}\|_{L^2(dx_j)}^4\,dx_{3-j}\leq Ch^3.
\tag{5.3}\label{eq:local-fiber-energy}
\]
Suppose that, for at least one $j$,
\[
\int_{\R}\mathcal U_{K/h}
\bigl(v_{j,m}(\cdot,x_{3-j})\bigr)\,dx_{3-j}
\leq C\rho h^3\quad\text{for every }m.
\tag{5.4}\label{eq:local-uniformity}
\]
Then
\[
\left\|T_{\mathrm{loc}}\left(\sum_m v_{1,m},\sum_n v_{2,n}\right)\right\|_1
\leq C\left[B_1B_2L^{1/2}h+(\rho+h+K^{-1})^{1/4}\right].
\tag{5.5}\label{eq:local-correlation-bound}
\]
Only indices whose cubes meet a fixed bounded region can contribute, so
the relevant sums are finite.
\end{lemma}

\begin{proof}
\textit{Spatial geometry.}
Call $(m,n)$ compatible when its contribution can be nonzero.
Compatibility implies
\[
hn=hm+\Gamma(t)+O(h)\quad\text{for some }t\in I.
\tag{5.6}\label{eq:compatible-pairs}
\]
There are $O(h^{-2})$ relevant indices in each family.

For fixed $m$, a $t$-interval of length $h$ permits only $O(1)$ values of
$n$, because $\Gamma$ is Lipschitz on $I$. The same assertion holds with
$m,n$ interchanged. Hence the number of compatible pairs is $O(h^{-3})$.

For every compatible pair choose a reference point with parameter
$\bar t_{mn}\in I$, and put
\[
q_{mn}=q(\bar t_{mn}).
\]
The complete $(x,y,t)$-support of that pair lies in a cube of side $Ch$.
In particular, its $(x,y)$-projection has area $O(h^2)$.

\medskip
\noindent\textit{Squaring and replacing $Q$ by its tangent.}
Write
\[
T_{mn}=T_{\mathrm{loc}}(v_{1,m},v_{2,n}).
\]
Insert a smooth cutoff equal to one on the pair's support and supported
in a cube of side $Ch$. Its derivatives of order $k$ are $O(h^{-k})$.

Cauchy--Schwarz in $(x,y)$, followed by writing the second time variable
as $t+s$, expresses $\|T_{mn}\|_1^2$ as at most $Ch^2$ times the absolute
value of an integral of
\[
D_s^{(1)}v_{1,m}(x+t,y)
\]
times
\[
v_{2,n}(x,y+Q(t+s))\overline{v_{2,n}(x,y+Q(t))}
\]
against a smooth amplitude supported where
\[
|s|+|t-\bar t_{mn}|+|x-\bar x|+|y-\bar y|\leq Ch.
\]
Here and below,
\[
D_s^{(j)}v(z)=v(z+se_j)\overline{v(z)}.
\]
Taylor's formula gives
\[
Q(t+s)-Q(t)=q_{mn}s+O(h^2)
\]
and
\[
Q(t)=Q(\bar t_{mn})+q_{mn}(t-\bar t_{mn})+O(h^2).
\]
Using \eqref{eq:local-derivatives}, the product can therefore be replaced by
\[
D_{q_{mn}s}^{(2)}v_{2,n}
\bigl(x,y+Q(\bar t_{mn})+q_{mn}(t-\bar t_{mn})\bigr)
\]
with pointwise error at most $CB_1^2B_2^2Lh^2$ after multiplication by the
first factor.

The four-dimensional integration region has measure $O(h^4)$. Including
the outside factor $h^2$, the error in $\|T_{mn}\|_1^2$ is at most
\[
CB_1^2B_2^2Lh^8.
\]
After taking square roots and summing over $O(h^{-3})$ pairs, its
contribution is
\[
CB_1B_2L^{1/2}h.
\tag{5.7}\label{eq:tangent-error}
\]

\medskip
\noindent\textit{Continuous Fourier expansion.}
Define
\[
a_{j,m}(s,k)=h^{-2}\widehat{D_s^{(j)}v_{j,m}}(k/h),
\qquad k\in\R^2.
\tag{5.8}\label{eq:correlation-fourier}
\]
Thus the inverse Fourier expansion uses the phase $e^{2\pi i k\cdot z/h}$.

For a fixed $s$, the Fourier kernel of the remaining integral is bounded by
\[
Ch^3\bigl(1+|k+\ell|+|k_1+q_{mn}\ell_2|\bigr)^{-12}.
\tag{5.9}\label{eq:fourier-kernel}
\]
Indeed, after replacing $Q$ by its tangent, the phase is linear in
$(x,y,t)$. Rescale the support to a fixed cube and integrate by parts
using $(1-\Delta)^6$. Assumption~\eqref{eq:smooth-cutoff-normalization}
supplies the required twelve derivatives.

Set
\[
w(u)=(1+|u|)^{-10}.
\]
Integrating the kernel in $\ell$, and using
\[
k_1+q\ell_2=(k_1-qk_2)+q(k_2+\ell_2),
\]
gives a row bound $Ch^3w(k_1-qk_2)$. The analogous column bound holds in
$k$. Cauchy--Schwarz with the nonnegative kernel therefore bounds the
remaining integral by
\[
Ch^3
\left(\int|a_{1,m}(s,k)|^2w(k_1-q_{mn}k_2)\,dk\right)^{1/2}
\left(\int|a_{2,n}(q_{mn}s,k)|^2w(k_1-q_{mn}k_2)\,dk\right)^{1/2}.
\tag{5.10}\label{eq:weighted-correlation}
\]
The Fourier manipulations are valid first after truncating the Fourier
variables. Passage to the limit follows from Plancherel: the original
bilinear integral is continuous in the two $L^2$ functions by
Cauchy--Schwarz at each $t$, while \eqref{eq:fourier-kernel} gives a bounded
$L^2$ kernel form.

Apply Cauchy--Schwarz in $s$, change variables $q_{mn}s\mapsto s$ in the
second factor, and then apply H\"older to the pair sum with exponents
$4,4,2$. Because there are $O(h^{-3})$ pairs, the remaining contribution
is at most
\[
Ch(H_1H_2)^{1/4},
\tag{5.11}\label{eq:correlation-sum}
\]
where
\[
H_j=\sum_{(m,n)\ \mathrm{compatible}}
\int_{\R}\int_{\R^2}|a_{j,m_j}(s,k)|^2
w(k_1-q_{mn}k_2)\,dk\,ds,
\]
with $m_1=m$ and $m_2=n$.

\medskip
\noindent\textit{Counting the possible slopes.}
Put
\[
G_j=\sum_m\int\!\!\int|a_{j,m}(s,k)|^2\,dk\,ds.
\]
Plancherel, \eqref{eq:autocorrelation-total}, and
\eqref{eq:local-fiber-energy} imply
\[
G_j=h^{-2}\sum_m\int\|v_{j,m}\|_{L^2(dx_j)}^4\,dx_{3-j}
\leq Ch^{-1}.
\tag{5.12}\label{eq:total-fourier-energy}
\]
For fixed $m_j$, the reference parameters of compatible partners have
bounded multiplicity in every interval of length $h$. Since $q'$ is
bounded below, for $k_2\neq0$,
\[
\sum_{\text{partners}}w(k_1-q_{mn}k_2)
\leq C\bigl(1+h^{-1}|k_2|^{-1}\bigr).
\tag{5.13}\label{eq:slope-count}
\]
To verify this, the parameters satisfying
\[
|k_1-q(t)k_2|\leq U
\]
lie in an interval of length at most $CU/|k_2|$. Count the reference
parameters in this interval and sum over $U=2^d$, $d\geq0$, with weights
$2^{-10d}$.

For $j=2$, \eqref{eq:slope-count} applies directly when $|k_2|>K$. For
$j=1$, if $|k_1|>K$, either $|k_2|\geq c|k_1|$, when
\eqref{eq:slope-count} applies, or
\[
|k_1-q(t)k_2|\geq c|k_1|
\]
throughout $I$. In the latter case the trivial count $O(h^{-1})$ and the
decay of $w$ give the same bound. Thus, for either $j$,
\[
\sum_{\text{partners}}w(k_1-q_{mn}k_2)
\leq C(1+h^{-1}K^{-1})\quad\text{when }|k_j|>K.
\tag{5.14}\label{eq:high-slope-count}
\]
For $|k_j|\leq K$, the trivial bound is $Ch^{-1}$.

Let
\[
L_j=\sum_m\int\!\!\int_{|k_j|\leq K}|a_{j,m}(s,k)|^2\,dk\,ds.
\]
For the family satisfying \eqref{eq:local-uniformity}, partial Plancherel
gives
\[
L_j\leq C\rho h^{-1}.
\tag{5.15}\label{eq:low-fourier-energy}
\]
Combining \eqref{eq:total-fourier-energy}--\eqref{eq:low-fourier-energy},
\[
H_j\leq Ch^{-2}(\rho+h+K^{-1})
\]
for that family. For the other family, the trivial partner count gives
\[
H_{3-j}\leq Ch^{-2}.
\]
Substitute these bounds in \eqref{eq:correlation-sum}, and add
\eqref{eq:tangent-error}. This proves \eqref{eq:local-correlation-bound}.
\end{proof}

\section{Local smoothing for bounded, frequency-localized inputs}

We now prove the local estimate from which the $L^2$ theorem will follow.

\begin{proposition}\label{prop:local-smoothing}
Fix $0<c_1<C_1<\infty$. Suppose
\[
\supp\widehat f_j\subseteq\{|\xi_j|\leq C_1L\},\qquad j=1,2,
\]
and, for at least one $j$,
\[
\supp\widehat f_j\subseteq\{c_1L\leq|\xi_j|\leq C_1L\}.
\]
Then, for $L\geq1$,
\[
\|T_{\mathrm{loc}}(f_1,f_2)\|_1
\leq C\|\psi\|_{C^{20}}L^{-\sigma}
\|f_1\|_\infty\|f_2\|_\infty,
\qquad\sigma=\frac{\kappa}{256}.
\tag{6.1}\label{eq:local-smoothing}
\]
\end{proposition}

\begin{proof}
\textit{Localization preserving the frequency restriction.}
Normalize the two $L^\infty$ norms and $\|\psi\|_{C^{20}}$ to be at most
one; zero inputs or cutoff give an immediate conclusion.

Choose smooth product cutoffs $\eta_m,\chi_m$, at spatial scale $h$, such
that
\[
\sum_m\eta_m=1,\qquad0\leq\chi_m\leq1,\qquad
\chi_m=1\text{ on }\supp\eta_m,
\]
with uniformly bounded overlap and derivative bounds $Ch^{-k}$. They can
be constructed from the partition used in Lemma~\ref{lem:decomposition}.

For each $j$, choose a smooth multiplier $P_j$, acting in coordinate $j$,
equal to one on the Fourier support of $f_j$. Its support is a slightly
larger ball, or a slightly larger annulus for the designated high-frequency
input. Put
\[
u_{j,m}=P_j(\eta_mf_j),\qquad
\widetilde f_j=\sum_m\chi_mu_{j,m}.
\tag{6.2}\label{eq:spatial-localization}
\]
The commutator identity
\[
\eta_mP_jf_j-P_j(\eta_mf_j)
=\int K_{j,L}(s)\bigl(\eta_m(z)-\eta_m(z-se_j)\bigr)
f_j(z-se_j)\,ds
\]
and \eqref{eq:kernel-norms} give
\[
\|f_j-\widetilde f_j\|_\infty\leq C(Lh)^{-1}.
\tag{6.3}\label{eq:commutator-error}
\]
Only boundedly many $\chi_m$ are nonzero at a given point. Also,
\[
\|\widetilde f_j\|_\infty\leq C.
\tag{6.4}\label{eq:localized-boundedness}
\]
Each $u_{j,m}$ has passive-coordinate support of length $Ch$, and
\[
\|u_{j,m}\|_\infty\leq C,\qquad
\|u_{j,m}\|_{L^2(dx_j)}^2\leq Ch.
\tag{6.5}\label{eq:localized-input-bounds}
\]
Apply Lemma~\ref{lem:decomposition} on each active-coordinate fiber, with
parameters
\[
R=K/h,\qquad\rho.
\]
Define
\[
v_{j,m,\sharp}=\chi_mu_{j,m,\sharp},\qquad
v_{j,m,\flat}=\chi_mu_{j,m,\flat}.
\]
Then
\[
\widetilde f_j=f_{j,\sharp}+f_{j,\flat},
\]
where each term is the sum over $m$.

By \eqref{eq:decomp-domination}, \eqref{eq:localized-flat-correlation}, and
\eqref{eq:localized-input-bounds},
\[
\int\|v_{j,m,\sharp}\|_{L^2(dx_j)}^4\,dx_{3-j}
+\int\|v_{j,m,\flat}\|_{L^2(dx_j)}^4\,dx_{3-j}\leq Ch^3,
\tag{6.6}\label{eq:decomposed-fiber-energy}
\]
and
\[
\int\mathcal U_R\bigl(v_{j,m,\flat}(\cdot,x_{3-j})\bigr)\,dx_{3-j}
\leq C\rho h^3.
\tag{6.7}\label{eq:decomposed-uniformity}
\]
Assume $R\leq cL$, with $c>0$ sufficiently small depending only on the
fixed frequency annulus. All nonzero frequency centers have magnitude at
most $CL$, and those for the designated high-frequency input have
magnitude at least $cL$.

The component bounds in Lemma~\ref{lem:decomposition} consequently imply
\[
\|v_{j,m,\sharp}\|_\infty+\|v_{j,m,\flat}\|_\infty\leq C\rho^{-1},
\]
and
\[
\|\partial_jv_{j,m,\sharp}\|_\infty
+\|\partial_jv_{j,m,\flat}\|_\infty\leq C\rho^{-1}L.
\tag{6.8}\label{eq:decomposed-derivatives}
\]

\medskip
\noindent\textit{The terms containing a flat part.}
By bilinearity, $T_{\mathrm{loc}}(\widetilde f_1,\widetilde f_2)$ is the sum
of
\[
T_{\mathrm{loc}}(f_{1,\flat},\widetilde f_2),\qquad
T_{\mathrm{loc}}(f_{1,\sharp},f_{2,\flat}),\qquad
T_{\mathrm{loc}}(f_{1,\sharp},f_{2,\sharp}).
\]
Lemma~\ref{lem:local-correlation}, together with
\eqref{eq:decomposed-fiber-energy}--\eqref{eq:decomposed-derivatives},
bounds the first two terms by
\[
C\left[\rho^{-2}L^{1/2}h+(\rho+h+K^{-1})^{1/4}\right].
\tag{6.9}\label{eq:flat-terms}
\]
The replacement errors in \eqref{eq:commutator-error} contribute at most
\[
C(Lh)^{-1},
\tag{6.10}\label{eq:operator-localization-error}
\]
using \eqref{eq:localized-boundedness}, boundedness of the original inputs,
and the fixed compact support of $\chi(x,y)\psi(t)$.

\medskip
\noindent\textit{An elementary nonstationary estimate.}
If $H\in C_c^1(J)$, $\phi\in C^2(J)$, and $\phi'$ is monotone, then, for
every $S>0$,
\[
\left|\int e^{2\pi i\phi(t)}H(t)\,dt\right|
\leq\int_{|\phi'(t)|\leq2S}|H(t)|\,dt
+\frac{C}{S}\bigl(\|H'\|_1+\|H\|_\infty\bigr).
\tag{6.11}\label{eq:nonstationary}
\]
To prove this, choose a smooth cutoff $\zeta$ equal to one on $[-1,1]$
and supported in $[-2,2]$. Separate the factor $\zeta(\phi'/S)$. For the
complement, integrate by parts using
\[
b(u)=\frac{1-\zeta(u/S)}{u},
\]
with $b(0)=0$. One has
\[
\|b\|_\infty\leq CS^{-1},\qquad
\int_{\R}|b'(u)|\,du\leq CS^{-1}.
\]
Monotonicity of $\phi'$ implies
\[
\int_J|b'(\phi'(t))\phi''(t)|\,dt\leq CS^{-1}.
\]
The integration-by-parts formula now gives \eqref{eq:nonstationary}.

\medskip
\noindent\textit{The sharp--sharp term.}
There are at most $C\rho^{-2}$ pairs of component indices. For each such
pair, the contribution is bounded by a sum over compatible spatial cubes
of integrals
\[
\iint\left|\int
 e^{2\pi i(\omega_{1,m}(y)t+\omega_{2,n}(x)Q(t))}
 H_{mn}(x,y,t)\,dt\right|\,dx\,dy.
\tag{6.12}\label{eq:sharp-oscillatory}
\]
The factors independent of $t$ have been removed inside the absolute value.

The amplitudes satisfy
\[
\|H_{mn}(x,y,\cdot)\|_\infty\leq C,\qquad
\|\partial_tH_{mn}(x,y,\cdot)\|_1\leq CRh.
\tag{6.13}\label{eq:sharp-amplitude}
\]
Only derivatives in the active coordinate are taken here. The frequency
functions are merely measurable in their passive coordinates, but those
coordinates remain fixed during differentiation in $t$.

The derivative of the phase is
\[
\omega_{1,m}(y)+q(t)\omega_{2,n}(x),
\]
which is monotone in $t$. Since $Rh=K\geq1$, \eqref{eq:nonstationary}, the
$O(h^2)$ spatial projection of each pair, and the $O(h^{-3})$ pair count
show that the nonstationary contributions total at most
\[
CR/S.
\tag{6.14}\label{eq:sharp-nonstationary}
\]
It remains to estimate the stationary contributions. Partition each
family of spatial cubes into finitely many classes so that cubes within
a class are disjoint. For example, using sufficiently large fixed residue
classes of $m\in\Z^2$ accomplishes this.

For fixed classes and fixed component indices, define measurable
functions on the input planes by assigning
\[
\alpha(u,y)=L^{-1}\omega_{1,m}(y)
\]
on the first-family cube indexed by $m$, and
\[
\beta(x,v)=-L^{-1}\omega_{2,n}(x)
\]
on the second-family cube indexed by $n$.

There is at most one possible index at each point. For the designated
high-frequency input, the assigned function has absolute value bounded
below by a fixed positive constant. Extend it outside the cubes with a
fixed nonzero value. Zero padded components may likewise be assigned a
high-frequency center without changing their value.

The stationary condition is now
\[
|\alpha(x+t,y)-q(t)\beta(x,y+Q(t))|\leq2S/L.
\]
At each $(x,y,t)$, disjointness permits at most one pair of cubes.
Lemma~\ref{lem:sublevel} therefore bounds the stationary contributions by
\[
C(S/L)^\kappa.
\tag{6.15}\label{eq:sharp-stationary}
\]
Combining \eqref{eq:sharp-nonstationary}, \eqref{eq:sharp-stationary}, and
the component count,
\[
\|T_{\mathrm{loc}}(f_{1,\sharp},f_{2,\sharp})\|_1
\leq C\rho^{-2}\bigl(R/S+(S/L)^\kappa\bigr).
\tag{6.16}\label{eq:sharp-terms}
\]

\medskip
\noindent\textit{Choosing all parameters.}
The estimates obtained so far give
\[
\|T_{\mathrm{loc}}(f_1,f_2)\|_1
\leq C\left[
(Lh)^{-1}+\rho^{-2}L^{1/2}h
+(\rho+h+K^{-1})^{1/4}
+\rho^{-2}\bigl(R/S+(S/L)^\kappa\bigr)\right].
\tag{6.17}\label{eq:all-local-errors}
\]
Choose
\[
h=L^{-3/4},\qquad K=L^{1/16},\qquad R=K/h=L^{13/16},
\]
\[
S=L^{7/8},\qquad\rho=L^{-\kappa/64}.
\tag{6.18}\label{eq:parameters}
\]
For sufficiently large $L$, all previous restrictions hold. The powers
occurring on the right of \eqref{eq:all-local-errors} are bounded by a
constant times the sum of
\[
L^{-1/4},\qquad L^{-1/4+\kappa/32},\qquad
L^{-\kappa/256},\qquad L^{-3/16},
\]
and
\[
L^{-1/64},\qquad L^{-1/16+\kappa/32},\qquad L^{-3\kappa/32}.
\]
Since $\kappa=1/39$, every term is at most $CL^{-\kappa/256}$.

The remaining bounded range of $L$ follows by the trivial localized
$L^\infty\times L^\infty\to L^1$ estimate, enlarging the constant.
Restoring the normalized factors proves Proposition~\ref{prop:local-smoothing}.
\end{proof}

\section{From the local bounded-input estimate to global \texorpdfstring{$L^2$}{L2} smoothing}

\subsection{Local \texorpdfstring{$L^2$}{L2} interpolation}

Let $P_1,P_2$ be fixed smooth frequency cutoffs at scale $L$, with one
annular and the other supported in a ball. Define
\[
S_L(u,v)=\chi T_\psi(P_1u,P_2v).
\]
Proposition~\ref{prop:local-smoothing} and \eqref{eq:young-l1} give
\[
\|S_L(u,v)\|_1
\leq C\|\psi\|_{C^{20}}L^{-\sigma}\|u\|_\infty\|v\|_\infty.
\tag{7.1}\label{eq:interpolation-upper}
\]
Equation~\eqref{eq:improving} and \eqref{eq:young-l1} give
\[
\|S_L(u,v)\|_1
\leq C\|\psi\|_{C^{20}}\|u\|_{12/7}\|v\|_{12/7}.
\tag{7.2}\label{eq:interpolation-lower}
\]
We prove the required interpolation directly. Set
\[
p_0=12/7,\qquad\theta=1-p_0/2=1/7.
\]
Normalize $\|u\|_2=\|v\|_2=1$, initially with $u,v$ simple and compactly
supported. On their nonzero sets define
\[
u_z=\frac{u}{|u|}|u|^{(2/p_0)(1-z)},\qquad
v_z=\frac{v}{|v|}|v|^{(2/p_0)(1-z)},
\]
and set these functions to zero elsewhere. For any $w\in L^\infty$ of
norm at most one, the function
\[
F(z)=\int S_L(u_z,v_z)w
\]
is analytic and bounded on the strip $0\leq\operatorname{Re}z\leq1$.
On its two boundaries,
\[
|F(it)|\leq C\|\psi\|_{C^{20}},
\]
\[
|F(1+it)|\leq C\|\psi\|_{C^{20}}L^{-\sigma}.
\]
The three-lines bound used here follows from the maximum-modulus
principle: divide $F(z)$ by $M_0^{1-z}M_1^z$, multiply by
$e^{\varepsilon z^2-\varepsilon}$, apply the maximum-modulus principle on
sufficiently tall rectangles, and let $\varepsilon\downarrow0$. The
horizontal sides tend to zero because of
$e^{-\varepsilon(\operatorname{Im}z)^2}$, while the vertical sides are
bounded by one. The maximum-modulus principle itself is available in
mathlib~\cite{mathlib-maximum}.

Since $u_\theta=u$ and $v_\theta=v$, taking the supremum over $w$ proves
\[
\|\chi T_\psi(P_1u,P_2v)\|_1
\leq C\|\psi\|_{C^{20}}L^{-\sigma/7}\|u\|_2\|v\|_2.
\tag{7.3}\label{eq:local-l2}
\]
Density and \eqref{eq:trivial} extend the estimate to arbitrary $L^2$ inputs.

\subsection{Globalization}

Let $Q_p=p+[0,1)^2$, $p\in\Z^2$, and let $\chi_q$ be translates of a
fixed smooth function equal to one on $Q_q$. Translation invariance makes
the constant in \eqref{eq:local-l2} uniform in $q$.

Suppose $f_j=P_jf_j$, with the prescribed ball or annular Fourier supports.
Write
\[
f_j=\sum_{p\in\Z^2}f_j\ind_{Q_p}.
\]
For any $N$, the kernel bound \eqref{eq:kernel-decay} gives
\[
\left\|\ind_{q+[-C,C]^2}P_j(f_j\ind_{Q_p})\right\|_2
\leq C_N(1+|p-q|)^{-N}\|f_j\ind_{Q_p}\|_2.
\tag{7.4}\label{eq:spatial-tail}
\]
To see this, the operator does not alter the passive coordinate, so
sufficiently separated passive intervals give zero. For separated active
intervals, \eqref{eq:kernel-decay}, followed by Cauchy--Schwarz on the unit
input interval, gives the stated decay. For nearby cubes, use the $L^2$
multiplier bound.

The elementary bound \eqref{eq:trivial}, localized to the region containing
all evaluations contributing to $\chi_qT_\psi$, gives
\[
\|\chi_qT_\psi(P_1(f_1\ind_{Q_p}),P_2(f_2\ind_{Q_s}))\|_1
\leq C\|\psi\|_{C^{20}}
(1+|p-q|)^{-8}(1+|s-q|)^{-8}
\|f_1\ind_{Q_p}\|_2\|f_2\ind_{Q_s}\|_2.
\tag{7.5}\label{eq:globalization-trivial}
\]
The same expression satisfies \eqref{eq:local-l2}, without the
spatial-decay factors. Taking the geometric mean of the two bounds yields
\[
C\|\psi\|_{C^{20}}L^{-\sigma/14}
(1+|p-q|)^{-4}(1+|s-q|)^{-4}
\|f_1\ind_{Q_p}\|_2\|f_2\ind_{Q_s}\|_2.
\tag{7.6}\label{eq:globalization-geometric}
\]
The sequence $(1+|p|)^{-4}$ belongs to $\ell^1(\Z^2)$. Sum
\eqref{eq:globalization-geometric} in $p,s,q$, apply Cauchy--Schwarz in $q$,
and use the discrete version of \eqref{eq:young-l1}. Since
\[
\sum_p\|f_j\ind_{Q_p}\|_2^2=\|f_j\|_2^2,
\]
we obtain
\[
\|T_\psi(f_1,f_2)\|_1
\leq C\|\psi\|_{C^{20}}L^{-s}\|f_1\|_2\|f_2\|_2,
\qquad s=\frac{\sigma}{14}=\frac1{139\,776}.
\tag{7.7}\label{eq:global-annular}
\]
Finite spatial sums suffice initially; passage to the full sums follows
from $L^2$ convergence and \eqref{eq:trivial}.

\section{The high-frequency tail, the \texorpdfstring{$C^1$}{C1} cutoff, and the original monomials}

\subsection{Removing the upper frequency restriction}

Choose $\varphi\in C_c^\infty(\R)$, equal to one on $[-1,1]$ and supported
in $[-2,2]$. Let $S_k^{(j)}$ be the multiplier
$\varphi(2^{-k}\xi_j)$, and put
\[
\Delta_k^{(j)}=S_k^{(j)}-S_{k-1}^{(j)}.
\]
The telescoping identity is
\[
T_\psi(f_1,f_2)
=\sum_{k\in\Z}T_\psi(\Delta_k^{(1)}f_1,S_k^{(2)}f_2)
+\sum_{k\in\Z}T_\psi(S_{k-1}^{(1)}f_1,\Delta_k^{(2)}f_2).
\tag{8.1}\label{eq:telescoping}
\]
Suppose one input has active Fourier support in
$\{|\xi_j|\geq\Lambda\}$, with $\Lambda\geq1$. Both summands vanish unless
$2^k\geq c\Lambda$. Each surviving term has one annular input at scale
$2^k$, while the other has active frequencies bounded by $C2^k$.

Apply \eqref{eq:global-annular}, use the uniform $L^2$ multiplier bounds,
and sum the geometric series. The finitely many possible scales below
one are covered by \eqref{eq:trivial}. This proves
\[
\|T_\psi(f_1,f_2)\|_1
\leq C\|\psi\|_{C^{20}}\Lambda^{-s}\|f_1\|_2\|f_2\|_2.
\tag{8.2}\label{eq:smooth-high-tail}
\]
The identity \eqref{eq:telescoping} follows first for finite telescoping
sums. Its upper endpoint converges in $L^1$ by \eqref{eq:trivial} and
$L^2$ convergence of $S_k^{(j)}f_j$. Its lower endpoint is zero because
of the high-frequency hypothesis.

\subsection{Reducing the cutoff requirement to \texorpdfstring{$C^1$}{C1}}

Now let $\psi\in C^1(\R)$ be supported in $I$. Choose a fixed nonnegative
smooth mollifier $\zeta$, supported in $[-1,1]$, with integral one. Put
\[
\psi_h=\psi*\zeta_h.
\]
For sufficiently small $h$, all $\psi_h$ are supported in one fixed
compact interval contained in $(0,\infty)$.

Differentiation of the mollifier gives
\[
\|\psi_h\|_{C^{20}}\leq Ch^{-20}\|\psi\|_\infty.
\tag{8.3}\label{eq:mollified-cutoff}
\]
The fundamental theorem of calculus and Fubini give
\[
\|\psi-\psi_h\|_1\leq Ch\|\psi'\|_1.
\tag{8.4}\label{eq:mollification-error}
\]
Apply \eqref{eq:smooth-high-tail} to $\psi_h$ and \eqref{eq:trivial} to
$\psi-\psi_h$. We obtain
\[
\|T_\psi(f_1,f_2)\|_1
\leq C\left(h^{-20}\Lambda^{-s}\|\psi\|_\infty
+h\|\psi'\|_1\right)\|f_1\|_2\|f_2\|_2.
\]
Choose
\[
h=c\Lambda^{-s/21},
\]
where $c>0$ is a fixed sufficiently small constant determined by $I$.
Then
\[
\|T_\psi(f_1,f_2)\|_1
\leq C\bigl(\|\psi\|_\infty+\|\psi'\|_1\bigr)
\Lambda^{-s/21}\|f_1\|_2\|f_2\|_2.
\tag{8.5}\label{eq:c1-smoothing}
\]
Thus the normalized curve $(t,t^r)$ satisfies the required theorem with
\[
\delta=\frac{s}{21}
=\frac{1}{39\cdot256\cdot14\cdot21}
=\frac{1}{2\,935\,296}.
\tag{8.6}\label{eq:final-exponent}
\]

\subsection{Returning to arbitrary distinct monomials}

Interchanging the two coordinates and the two inputs preserves all norms
and the active-frequency hypothesis. We may therefore assume
\[
\alpha_1=m<n=\alpha_2.
\]
Set
\[
u=t^m,\qquad r=n/m>1.
\]
The original operator becomes
\[
\int_{a^m}^{b^m}\widetilde\psi(u)
 g_1(x+u,y)g_2(x,y+u^r)\,du,
\]
where
\[
\widetilde\psi(u)=\frac1m u^{1/m-1}\psi(u^{1/m}).
\tag{8.7}\label{eq:changed-cutoff}
\]
On the fixed interval $[a^m,b^m]$, the weight and its derivative are
bounded. The chain rule and the substitution $u=t^m$ give
\[
\|\widetilde\psi\|_\infty+\|\widetilde\psi'\|_1
\leq C_{a,b,m}\bigl(\|\psi\|_\infty+\|\psi'\|_1\bigr).
\tag{8.8}\label{eq:changed-cutoff-bound}
\]
No spatial variable is changed, so the Fourier-support hypothesis and the
input $L^2$ norms are unchanged. Applying \eqref{eq:c1-smoothing} proves
\eqref{eq:target}.

In the original parameter coordinates, the modified flow is therefore
described by
\[
\theta_1(\mathbf t)=-t_1^m+t_2^m-t_3^m,\qquad
\theta_2(\mathbf t)=t_1^n-t_2^n+t_3^n,
\]
\[
h(t)=\frac nm t^{n-m},\qquad
\nu(\mathbf t)=\frac{h(t_2)}{h(t_1)h(t_3)},
\]
and
\[
V=\nabla\theta_1\times\nabla\theta_2.
\]
The coordinate substitution $u_i=t_i^m$ converts this field, up to a
positive time-rescaling factor, into \eqref{eq:vector-field}.

The distinct-exponent assumption is used quantitatively in the positive
lower bound for $q'$ and in the determinant estimate
\eqref{eq:flow-nondegeneracy}. The cutoff dependence is handled separately
by \eqref{eq:mollified-cutoff}--\eqref{eq:c1-smoothing}, so neither the decay
exponent nor the implicit constant needs any additional dependence on the
particular cutoff $\psi$ beyond the displayed factor
$\|\psi\|_\infty+\|\psi'\|_1$.

\begin{thebibliography}{9}
\raggedright

\bibitem{CDR}
M.~Christ, P.~Durcik, and J.~Roos,
\emph{arXiv:2008.10140}.
\url{https://arxiv.org/abs/2008.10140}.

\bibitem{mathlib-fourier}
The mathlib community,
\emph{Mathlib.Analysis.Fourier.LpSpace}.
\url{https://leanprover-community.github.io/mathlib4_docs/Mathlib/Analysis/Fourier/LpSpace.html}.

\bibitem{mathlib-maximum}
The mathlib community,
\emph{Mathlib.Analysis.Complex.AbsMax}.
\url{https://leanprover-community.github.io/mathlib4_docs/Mathlib/Analysis/Complex/AbsMax.html}.

\end{thebibliography}

\end{document}
